\textbf{Proof.}
The published comparator supplies only the normalization recorded in \(\mathrm{def:ryan-relative-efficiency}\); the following calculation is internal to the present finite menu.  By \(\mathrm{thm:quadruple-menu-frontier}\), \(r_j^\star=4c_j/n^2\).  Its trace certificates imply, with \(p_k=p_k(\pi)\),
\[
\operatorname{RE}_1(\pi)\le\frac1{1+F_1},
\qquad
\operatorname{RE}_2(\pi)\le\frac1{1+F_2}, \tag{1}
\]
where
\[
F_1=\kappa_1(ap_2+bp_3),
\qquad F_2=\kappa_2(ap_1+bp_3). \tag{2}
\]
All denominators are positive.  It therefore suffices to minimize \(\max(F_1,F_2)\), provided the trace bounds can be attained simultaneously.

Assume \(a>2b>0\) and \(\kappa_2>\kappa_1>0\), and put
\[
D=\kappa_1(a-b)+\kappa_2b,
\qquad
\alpha=\frac{\kappa_2b}{D},
\qquad
1-\alpha=\frac{\kappa_1(a-b)}D. \tag{3}
\]
Then \(D>0\) and \(0<\alpha<1\).  Substituting \(p_3=1-p_1-p_2\) gives the affine certificate
\[
\alpha F_1+(1-\alpha)F_2
=\frac{ab\kappa_1\kappa_2}{D}
+\frac{a(a-2b)\kappa_1\kappa_2}{D}p_1
\ge m:=\frac{ab\kappa_1\kappa_2}{D}. \tag{4}
\]
Since a maximum dominates every convex combination, \(\max(F_1,F_2)\ge m\).

Set
\[
p_1^\circ=0,
\qquad
p_2^\circ=\frac{b(\kappa_2-\kappa_1)}D,
\qquad
p_3^\circ=1-p_2^\circ. \tag{5}
\]
The second mass is positive, while
\[
D-b(\kappa_2-\kappa_1)=a\kappa_1>0, \tag{6}
\]
so (5) is in the simplex.  Substitution gives \(F_1=F_2=m\).  The independent-block, independent-sign law with (5) lies in \(\Pi_n\) and has scalar score matrices
\[
C_j=4c_j(1+F_j)I_q,
\qquad
R_j=\frac{4c_j}{n^2}(1+F_j). \tag{7}
\]
Thus (1) is attained for both classes and
\[
\eta_{\mathrm{REL}}
=\frac1{1+m}
=\left[1+\frac{ab\kappa_1\kappa_2}
{\kappa_1(a-b)+\kappa_2b}\right]^{-1}. \tag{8}
\]

If any design attains (8), then both inequalities
\[
\min_j\operatorname{RE}_j(\pi)
\le\{1+\max(F_1,F_2)\}^{-1}
\le(1+m)^{-1} \tag{9}
\]
are equalities.  The coefficient of \(p_1\) in (4) is strictly positive, so \(p_1=0\).  Since \(0<\alpha<1\), equality of the maximum and the convex combination forces \(F_1=F_2=m\).  Solving on the face \(p_1=0\) gives exactly (5).  By contrast, the \(a>2b\) branch of \(\mathrm{thm:quadruple-menu-frontier}\) forces every gain-normalized optimizer to have \(p_3=1\), whereas (5) has \(p_2^\circ>0\).  The optimizer sets are therefore disjoint.  Both calculations use the same oracle risks, committed risks, and feasible laws; only the criterion normalization differs, so no submodel claim about the comparator is needed. \(\square\)